% GENERATED FILE. Edit canonical lessons, then run npm run generate:books.
% canonical-source-hash: fnv1a64:de588ce4097808f6
% canonical-lessons: MR-C318-janma, MR-C318-barse, MR-C318-vadhdivas, MR-C318-munj, MR-C318-patrika

\chapter{A New Baby}
\label{ch:mr-a-new-baby}
\hlchaptermodality{318}

\begin{chapteropening}
\textbf{By the end of this chapter:} \emph{I can talk about births, naming days and birthdays.}

Complete the last lesson of chapter 318: I can talk about births, naming days and birthdays.
\end{chapteropening}

\section[\texorpdfstring{\mr{जन्म} (janma)}{जन्म (janma)}]{\mr{जन्म} (janma) — birth}

\label{lesson:MR-C318-janma}



\begin{quote}
Before the new one: say the Marathi for a shehnai, then the Marathi for a rhyming couplet.
\end{quote}

\subsection*{You'll want to know: \mr{जन्म}}
\textbf{\mr{जन्म}} — \emph{janma} — \textquotedblleft{}birth\textquotedblright{}.

\begin{quote}
\mr{माझ्या} \mr{भावाचा} \mr{जन्म} \mr{पुण्यात} \mr{झाला}.
\end{quote}

\begin{quote}
\emph{mājhyā bhāvācā janma puṇyāt jhālā.}
\end{quote}

\begin{quote}
My brother was born in Pune.
\end{quote}

Masculine.

X-\mr{चा} \mr{जन्म} \mr{होणे}, X is born.

\begin{cousinweb}
\textbf{\mr{जन्म}} (\emph{janma}) is a tatsama, taken from Sanskrit \emph{janman}, 'birth' (Monier-Williams).
\end{cousinweb}

\begin{grammarlens}[title={What you've built}]
One more word for telling what happened.
\end{grammarlens}

\subsection*{Your turn}
\begin{itemize}
\raggedright
  \item \emph{Say it:} \emph{janma}
  \item \emph{Say it:} \emph{janma}, once more
  \item \emph{Say it:} \emph{mājhyā bhāvācā janma puṇyāt jhālā.} — My brother was born in Pune.
  \item \emph{Recall:} say the Marathi for a shehnai, then the Marathi for a rhyming couplet, then say \emph{janma} again
\end{itemize}

\begin{tcolorbox}[breakable,colback=teal!4,colframe=teal!35!black,title={Before you move on}]
What does \mr{जन्म} mean? (\textquotedblleft{}Birth\textquotedblright{}.) Say it once more.
\end{tcolorbox}
\section[\texorpdfstring{\mr{बारसे} (bārse)}{बारसे (bārse)}]{\mr{बारसे} (bārse) — a naming ceremony}

\label{lesson:MR-C318-barse}



\begin{quote}
Before the new one: say the Marathi for a rhyming couplet, then the Marathi for birth.
\end{quote}

\subsection*{You'll want to know: \mr{बारसे}}
\textbf{\mr{बारसे}} — \emph{bārse} — \textquotedblleft{}a naming ceremony\textquotedblright{}.

\begin{quote}
\mr{बाळाचं} \mr{बारसं} \mr{बाराव्या} \mr{दिवशी} \mr{झालं}.
\end{quote}

\begin{quote}
\emph{bāḷācaṁ bārsaṁ bārāvyā divśī jhālaṁ.}
\end{quote}

\begin{quote}
The baby's naming ceremony was on the twelfth day.
\end{quote}

Neuter.

Neuter, written \mr{बारसं} in a sentence like \mr{मोठं}; held on the twelfth day after a birth, from \mr{बारा}, twelve.

\begin{cousinweb}
\textbf{\mr{बारसे}} (\emph{bārse}) is built on \emph{bārā}, 'twelve', the day the rite falls on; CDIAL lists Marathi \emph{bārāvā}, 'twelfth', under Sanskrit \emph{dvādaśama} (CDIAL 6660).
\end{cousinweb}

\begin{grammarlens}[title={What you've built}]
One more word for telling what happened.
\end{grammarlens}

\subsection*{Your turn}
\begin{itemize}
\raggedright
  \item \emph{Say it:} \emph{bārse}
  \item \emph{Say it:} \emph{bārse}, once more
  \item \emph{Say it:} \emph{bāḷācaṁ bārsaṁ bārāvyā divśī jhālaṁ.} — The baby's naming ceremony was on the twelfth day.
  \item \emph{Recall:} say the Marathi for a rhyming couplet, then the Marathi for birth, then say \emph{bārse} again
\end{itemize}

\begin{tcolorbox}[breakable,colback=teal!4,colframe=teal!35!black,title={Before you move on}]
What does \mr{बारसे} mean? (\textquotedblleft{}A naming ceremony\textquotedblright{}.) Say it once more.
\end{tcolorbox}
\section[\texorpdfstring{\mr{वाढदिवस} (vāḍhdivas)}{वाढदिवस (vāḍhdivas)}]{\mr{वाढदिवस} (vāḍhdivas) — a birthday}

\label{lesson:MR-C318-vadhdivas}



\begin{quote}
Before the new one: say the Marathi for birth, then the Marathi for a naming ceremony.
\end{quote}

\subsection*{You'll want to know: \mr{वाढदिवस}}
\textbf{\mr{वाढदिवस}} — \emph{vāḍhdivas} — \textquotedblleft{}a birthday\textquotedblright{}.

\begin{quote}
\mr{काल} \mr{माझ्या} \mr{मुलीचा} \mr{पाचवा} \mr{वाढदिवस} \mr{होता}.
\end{quote}

\begin{quote}
\emph{kāl mājhyā mulīcā pācvā vāḍhdivas hotā.}
\end{quote}

\begin{quote}
Yesterday was my daughter's fifth birthday.
\end{quote}

Masculine.

\mr{वाढदिवसाच्या} \mr{शुभेच्छा}, happy birthday.

\begin{cousinweb}
\textbf{\mr{वाढदिवस}} (\emph{vāḍhdivas}) is built from \emph{vāḍh}, 'increase, growth', which CDIAL lists under Sanskrit \emph{vardha} (CDIAL 11373), and \emph{divas}, 'a day': the day one grows a year.
\end{cousinweb}

\begin{grammarlens}[title={What you've built}]
One more word for telling what happened.
\end{grammarlens}

\subsection*{Your turn}
\begin{itemize}
\raggedright
  \item \emph{Say it:} \emph{vāḍhdivas}
  \item \emph{Say it:} \emph{vāḍhdivas}, once more
  \item \emph{Say it:} \emph{kāl mājhyā mulīcā pācvā vāḍhdivas hotā.} — Yesterday was my daughter's fifth birthday.
  \item \emph{Recall:} say the Marathi for birth, then the Marathi for a naming ceremony, then say \emph{vāḍhdivas} again
\end{itemize}

\begin{tcolorbox}[breakable,colback=teal!4,colframe=teal!35!black,title={Before you move on}]
What does \mr{वाढदिवस} mean? (\textquotedblleft{}A birthday\textquotedblright{}.) Say it once more.
\end{tcolorbox}
\section[\texorpdfstring{\mr{मुंज} (muñj)}{मुंज (muñj)}]{\mr{मुंज} (muñj) — a thread ceremony}

\label{lesson:MR-C318-munj}



\begin{quote}
Before the new one: say the Marathi for a naming ceremony, then the Marathi for a birthday.
\end{quote}

\subsection*{You'll want to know: \mr{मुंज}}
\textbf{\mr{मुंज}} — \emph{muñj} — \textquotedblleft{}a thread ceremony\textquotedblright{}.

\begin{quote}
\mr{माझ्या} \mr{भावाच्या} \mr{मुलाची} \mr{मुंज} \mr{पुढच्या} \mr{आठवड्यात} \mr{आहे}.
\end{quote}

\begin{quote}
\emph{mājhyā bhāvācyā mulācī muñj puḍhacyā āṭhvaḍyāt āhe.}
\end{quote}

\begin{quote}
My brother's son's thread ceremony is next week.
\end{quote}

Feminine.

The Marathi name for the upanayana, the rite at which a boy is given the sacred thread; it is kept in some communities.

\begin{cousinweb}
\textbf{\mr{मुंज}} (\emph{muñj}) is a tadbhava: CDIAL lists Marathi \emph{mũjī}, 'the sacred thread', under Sanskrit \emph{mauñjī}, the cord of \emph{muñja} grass worn by a Brahman (CDIAL 10367; \emph{muñja} is CDIAL 10184). The rite is named after that cord.
\end{cousinweb}

\begin{grammarlens}[title={What you've built}]
One more word for telling what happened.
\end{grammarlens}

\subsection*{Your turn}
\begin{itemize}
\raggedright
  \item \emph{Say it:} \emph{muñj}
  \item \emph{Say it:} \emph{muñj}, once more
  \item \emph{Say it:} \emph{mājhyā bhāvācyā mulācī muñj puḍhacyā āṭhvaḍyāt āhe.} — My brother's son's thread ceremony is next week.
  \item \emph{Recall:} say the Marathi for a naming ceremony, then the Marathi for a birthday, then say \emph{muñj} again
\end{itemize}

\begin{tcolorbox}[breakable,colback=teal!4,colframe=teal!35!black,title={Before you move on}]
What does \mr{मुंज} mean? (\textquotedblleft{}A thread ceremony\textquotedblright{}.) Say it once more.
\end{tcolorbox}
\section[\texorpdfstring{\mr{पत्रिका} (patrikā)}{पत्रिका (patrikā)}]{\mr{पत्रिका} (patrikā) — an invitation card}

\label{lesson:MR-C318-patrika}



\begin{quote}
Before the new one: say the Marathi for a birthday, then the Marathi for a thread ceremony.
\end{quote}

\subsection*{You'll want to know: \mr{पत्रिका}}
\textbf{\mr{पत्रिका}} — \emph{patrikā} — \textquotedblleft{}an invitation card\textquotedblright{}.

\begin{quote}
\mr{त्यांनी} \mr{लग्नाची} \mr{पत्रिका} \mr{घरी} \mr{आणून} \mr{दिली}.
\end{quote}

\begin{quote}
\emph{tyānnī lagnācī patrikā gharī āṇūn dilī.}
\end{quote}

\begin{quote}
They brought the wedding invitation card to our home.
\end{quote}

Feminine.

\mr{लग्नपत्रिका}, a wedding card; also a small journal.

\begin{cousinweb}
\textbf{\mr{पत्रिका}} (\emph{patrikā}) is built on Sanskrit \emph{patra}, 'a leaf, a letter' (Monier-Williams), with the small-thing ending \emph{-ikā}.
\end{cousinweb}

\begin{grammarlens}[title={What you've built}]
One more word for telling what happened.
\end{grammarlens}

\subsection*{Your turn}
\begin{itemize}
\raggedright
  \item \emph{Say it:} \emph{patrikā}
  \item \emph{Say it:} \emph{patrikā}, once more
  \item \emph{Say it:} \emph{tyānnī lagnācī patrikā gharī āṇūn dilī.} — They brought the wedding invitation card to our home.
  \item \emph{Recall:} say the Marathi for a birthday, then the Marathi for a thread ceremony, then say \emph{patrikā} again
\end{itemize}

\begin{tcolorbox}[breakable,colback=teal!4,colframe=teal!35!black,title={Before you move on}]
What does \mr{पत्रिका} mean? (\textquotedblleft{}An invitation card\textquotedblright{}.) Say it once more.
\end{tcolorbox}
